\documentclass{article}
\usepackage[UTF8]{ctex}
\usepackage{amsmath,amssymb,amsfonts}
\usepackage{enumitem}
\usepackage{geometry}
\geometry{a4paper, margin=1in}
% === 错题本排版优化（不参与题目标识） ===
\setlength{\parindent}{0pt}
\setlength{\parskip}{0.35em plus 0.10em minus 0.08em}
\setlength{\emergencystretch}{3em}
\setlist{itemsep=0.25em, topsep=0.40em, parsep=0pt}
\allowdisplaybreaks[2]
% === 错题本排版优化结束 ===
\begin{document}
\section*{安徽大学2024年高等代数真题}
\section*{103．安徽大学2024年高等代数试卷}
一．填空题（本大题共六小题，每小题5分，共30分）．

\begin{itemize}

\subsection{题目1}
  \item[1．] 设
\[
A=\left(\begin{array}{ll}
1 & 2 \\
2 & 3
\end{array}\right), B=\left(\begin{array}{ccc}
1 & 0 & 0 \\
-1 & 1 & 0 \\
0 & 2 & 1
\end{array}\right), C=\left(\begin{array}{ccc}
0 & 1 & 0 \\
-1 & 0 & 2
\end{array}\right) .
\]
若矩阵满足 \(A X B=C\) ，则 \(X=\) \(\_\_\_\_\)．

\subsection{题目2}
  \item[2．] 设 \(A=\left(a_{i j}\right)\) 为2阶复方阵，满足 \(\operatorname{tr}\left(A^{k}\right)=k(k=1,2)\),\(R ，其中 \operatorname{tr}(A)=a_{11}+a_{22}\) 为矩阵 \(A\) 的迹，则行列式 \(|A|=\) \(\_\_\_\_\)。
\end{itemize}

\subsection{题目3}
3．设矩阵 \(A=\left(\begin{array}{cccc}-2 & 1 & 1 & 1 \\ -1 & -1 & 0 & 1 \\ 0 & 0 & -1 & 0 \\ -1 & 1 & 1 & 0\end{array}\right)\) ，则 \(A\) 的最小多项式为 \(\_\_\_\_\)．\\
\(\_\_\_\_\)

\subsection{题目4}
4．设 \(n\) 阶矩阵 \(A\) 的各行元素之和均为 \(3, E\) 为单位阵，则矩阵 \(A^{2}-2 A+E\) 的各行元素之和为\\
\(\_\_\_\_\)。

\subsection{题目5}
5．设 \(\mathbb{R}^{3}\) 为带标准内积的3维欧氏空间，对 \(\mathbb{R}^{3}\) 的基 \(\alpha_{1}=(-1,1,1), \alpha_{2}=(0,-1,1), \alpha_{3}=(0,0,1)\) 所 Schmidt 正交化得 \(\mathbb{R}^{3}\) 的标准正交基 \(\beta_{1}, \beta_{2}, \beta_{3}\) ，则 \(\beta_{3}=\) \(\_\_\_\_\)．

\subsection{题目6}
6．设 \(\lambda-\) 矩阵

\[
A(\lambda)=\left(\begin{array}{lll}
(\lambda-1)^{3} & & \\
& (\lambda+1)^{2} & \\
& & (\lambda-1)^{2}(\lambda+1)^{4}
\end{array}\right) .
\]

则 \(A(\lambda)\) 的所有不变因子为 \(\_\_\_\_\)．

二．辨析题（本题共四小题，每小题5分，共20分．先判断对错，再简要说明）．\\

\subsection{题目7}
7．设 \(f(x)\) 为实系数多项式，且次数为奇数，则 \(f(x)\) 必有实根．

\subsection{题目8}
8．设 \(A, B\) 都为4阶复方阵，则 \(A\) 与 \(B\) 相似当且仅当 \(A\) 与 \(B\) 有相同的特征多项式，且每个特征值的几何重数（即对应特征子空间的维数）也相同。

\subsection{题目9}
9．设 \(\sigma\) 是 \(n\) 维复线性空间 \(V\) 上的线性变换，\(n>1\) ，若 \(\sigma^{n}=0\) 且 \(\sigma^{n-1} \neq 0\) ，则存在两个 \(\sigma\) 的非平凡子空间 \(U\) 和 \(W\) ，使得 \(V=U \oplus W\) 。

\subsection{题目10}
10．设 \(A(\lambda), B(\lambda)\) 都是数域 \(P\) 上 \(m \times n\) 的 \(\lambda\)－矩阵，则 \(A(\lambda), B(\lambda)\) 等价的充要条件为 \(A(\lambda)\) 与\(B(\lambda)\) 有相同的初等因子组．

三。计算题（本大题共五小题，每小题10分，共50分）。\\

\subsection{题目11}
11．设 \(A=\left(a_{i j}\right)\) ，其中 \(a_{i j}=\frac{1-x_{i}^{n} y_{j}^{n}}{1-x_{i} y_{j}}, x_{i} y_{j} \neq 1, i, j=1,2, \cdots, n\) ，求行列式 \(|A|\) 。

\subsection{题目12}
12．设多项式

\[
f(x)=2 x^{4}+2 x^{3}-5 x^{2}-13 x-2, g(x)=x^{3}+x^{2}-3 x-6 .
\]

求 \(f(x)\) 与 \(g(x)\) 的首一最大公因式（ \(f(x), g(x)\) ）以及多项式 \(u(x)\) 与 \(v(x)\) ，使得

\[
u(x) f(x)+v(x) g(x)=(f(x), g(x)) .
\]

\subsection{题目13}
13．设 \(A\) 是 \(5 \times 4\) 矩阵，且秩 \((A)=3, \beta\) 为5维非零列向量。已知 \(\gamma_{1}, \gamma_{2}, \gamma_{3}\) 为方程组 \(A X=\beta\) 的 3个不同的解，且 \(\gamma_{1}+\gamma_{2}=(2,2,0,2)^{\mathrm{T}}, \gamma_{1}+\gamma_{3}=(0,0,2,0)^{\mathrm{T}}\) ．求 \(A X=\beta\) 的通解．

\subsection{题目14}
14．设复二次型 \(f\left(x_{1}, x_{2}, x_{3}\right)=2 x_{1}^{2}+x_{2}^{2}-4 x_{1} x_{2}-4 x_{2} x_{3}\) ，求非退化线性替换 \(X=C Y\) ，将二次型\(f\left(x_{1}, x_{2}, x_{3}\right)\) 化为规范形，其中 \(X=\left(x_{1}, x_{2}, x_{3}\right)^{\mathrm{T}}, Y=\left(y_{1}, y_{2}, y_{3}\right)^{\mathrm{T}}\) ，并写出规范形．

\subsection{题目15}
15．设实矩阵

\[
A=\left(\begin{array}{lll}
-1 & 1 & 2 \\
-1 & 1 & 1 \\
-2 & 1 & 3
\end{array}\right)
\]

\begin{itemize}
  \item[（1）] 求 \(A\) 的若尔当标准形 \(J\) 。
  \item[（2）] 求可逆矩阵 \(P\) ，使得 \(P^{-1} A P=J\) 。
\end{itemize}

四．证明题（本大题共四小题，第16，17题每题10分，第18，19题每题15分，共50分）．\\

\subsection{题目16}
16．设多项式 \(f(x)=x^{p}+p x+p-1\) ，其中 \(p\) 为奇素数．证明：\(f(x)\) 在有理数域上不可约．

\subsection{题目17}
17．设 \(n\) 元实二次型 \(f\left(x_{1}, x_{2}, \cdots, x_{n}\right)=l_{1}^{2}+\cdots+l_{a}^{2}-l_{s+1}^{2}-\cdots-l_{s+t}^{2}\) ，其中

\[
l_{i}=a_{i 1} x_{1}+a_{i 2} x_{2}+\cdots+a_{i n} x_{n}, a_{i j} \in \mathbb{R}, i=1,2, \cdots, s+t, j=1,2, \cdots, n .
\]

证明：\(f\left(x_{1}, x_{2}, \cdots, x_{n}\right)\) 的正惯性指数 \(p \leq s\) 。

\subsection{题目18}
18．设 \(\sigma\) 为 \(n\) 维欧氏空间 \(V\) 上的线性变换，且在基 \(\alpha_{1}, \alpha_{2}, \cdots, \alpha_{n}\) 下的矩阵为 \(A\) 。证明：\(\sigma\) 为对称变换的充要条件是 \(A^{\mathrm{T}} G=G A\) ，其中 \(G=\left(\left(\alpha_{i}, \alpha_{j}\right)\right)\) 为基 \(\alpha_{1}, \alpha_{2}, \cdots, \alpha_{n}\) 的度量矩阵。

\subsection{题目19}
19．设 \(A\) 为 \(n\) 阶复方阵， 0 为 \(A\) 的最小多项式 \(m(\lambda)\) 的 \(r\) 重根，\(r \geq 2\) 为正整数。证明：

\begin{itemize}
  \item[（1）] 对任意的正整数 \(k \geq r\) ，秩 \(\left(A^{k}\right)=\) 秩 \(\left(A^{r}\right)\) 。
  \item[（2）] 秩 \(\left(A^{r}\right)<\) 秩 \(\left(A^{r-1}\right)\) 。
\end{itemize}
\end{document}
